% \usepackage{cmap} % Улучшенный поиск русских слов в полученном pdf-файле
% \usepackage[T2A]{fontenc} % Поддержка русских букв
% \usepackage[utf8]{inputenc} % Кодировка utf8
% \usepackage[english,russian]{babel} % Языки: русский, английский
% \usepackage{CJKutf8}

% \usepackage[14pt]{extsizes}

% \usepackage[table]{xcolor}

% \usepackage{graphicx}
% \usepackage{multirow}

% \usepackage{caption}
% \captionsetup{labelsep=endash}
% \captionsetup[figure]{name={Рисунок}}

% \usepackage{amsmath}
% \usepackage{amsfonts}

% \usepackage{geometry}
% \geometry{left=30mm}
% \geometry{right=10mm}
% \geometry{top=20mm}
% \geometry{bottom=20mm}

% \usepackage{alphabeta}
% \usepackage{enumitem}
% \AddEnumerateCounter{\asbuk}{\@asbuk}{м}

% \usepackage{tabularx}

% % Переопределение стандартных \section, \subsection, \subsubsection по ГОСТу;

% \usepackage{titlesec}[explicit]
% \titleformat{name=\section,numberless}[block]{\normalfont\large\hyphenpenalty=10000\bfseries\centering}{}{0pt}{}
% \titleformat{\section}[block]{\normalfont\large\bfseries}{\thesection}{1em}{}
% \titlespacing\section{\parindent}{*4}{*4}

% \titleformat{\subsection}[hang]
% {\bfseries\large}{\thesubsection}{1em}{}
% \titlespacing\subsection{\parindent}{*2}{*2}

% \titleformat{\subsubsection}[hang]
% {\bfseries\large}{\thesubsubsection}{1em}{}
% \titlespacing\subsubsection{\parindent}{*2}{*2}

% \usepackage[hyphens]{url}

% % Переопределение их отступов до и после для 1.5 интервала во всем документе
% \usepackage{setspace}
% \onehalfspacing % Полуторный интервал
% \frenchspacing
% \setlength\parindent{1.25cm}

% \usepackage{indentfirst} % Красная строка

% % Настройки оглавления
% \usepackage{multirow}
% \usepackage{longtable}
% \usepackage{etoolbox}
% \makeatletter
% \patchcmd{\tableofcontents}{\@starttoc{toc}}{\vspace{-1cm}\@starttoc{toc}}{}{}
% \makeatother
% \usepackage{tocloft}
% \renewcommand{\cftsecleader}{\cftdotfill{\cftdotsep}}
% \renewcommand{\cfttoctitlefont}{\hfill\large\bfseries}
% \renewcommand{\cftaftertoctitle}{\hfill\hfill}
% \setlength\cftbeforetoctitleskip{-22pt}
% \setlength\cftaftertoctitleskip{15pt}

% % Гиперссылки
% \usepackage[pdftex]{hyperref}
% % \usepackage[xetex]{hyperref}
% % \usepackage{hyperref}
% \hypersetup{hidelinks}

% % \numberwithin{equation}{section}

% % Дополнительное окружения для подписей
% \usepackage{array}
% \newenvironment{signstabular}[1][1]{
% 	\renewcommand*{\arraystretch}{#1}
% 	\tabular
% }{
% 	\endtabular
% }

% \usepackage{makecell}
% \usepackage{longtable}

% \usepackage{enumitem} 
% \setenumerate[0]{label=\arabic*)} % Изменение вида нумерации списков
% \renewcommand{\labelitemi}{---}

% % Листинги 
% \usepackage{courier}

% \usepackage{float}
% \usepackage{listings}
% \floatstyle{plaintop}
% \newfloat{code}{H}{myc}

% \usepackage{chngcntr} % Listings counter within section set main.tex after begin document

% % Для листинга кода:
% \lstset{
% 	basicstyle=\small\ttfamily,			% размер и начертание шрифта для подсветки кода
% 	language=SQL,   					% выбор языка для подсветки	
% 	numbers=left,						% где поставить нумерацию строк (слева\справа)
% 	numbersep=5pt,
% 	% stepnumber=1,						% размер шага между двумя номерами строк
% 	xleftmargin=17pt,
% 	% showstringspaces=false,
% 	numbersep=5pt,						% как далеко отстоят номера строк от подсвечиваемого кода
% 	frame=single,						% рисовать рамку вокруг кода
% 	tabsize=4,							% размер табуляции по умолчанию равен 4 пробелам
% 	captionpos=t,						% позиция заголовка вверху [t] или внизу [b]
% 	breaklines=true,					
% 	breakatwhitespace=true,				% переносить строки только если есть пробел
% 	escapeinside={\#*}{*)},				% если нужно добавить комментарии в коде
% 	inputencoding=utf8x,
% 	backgroundcolor=\color{white},
% 	numberstyle=,%\tiny,					    % размер шрифта для номеров строк
% 	% keywordstyle=\color{blue},
% 	keywordstyle=\color{black},
% 	% stringstyle=\color{red!90!black}, % color of text in ""
% 	stringstyle=\color{black}, % color of text in ""
% 	% commentstyle=\color{green!50!black},
% 	commentstyle=\color{black},
%     morekeywords={TO,WITH,SUPERUSER,PRIVILEGES,TABLES,SCHEMA,USER,PASSWORD,public,REFERENCES,UUID,IF}
% }

% \lstdefinelanguage{Golang}%
%   {morekeywords=[1]{package,import,func,type,struct,return,defer,panic,%
%      recover,select,var,const,iota,},%
%    morekeywords=[2]{string,uint,uint8,uint16,uint32,uint64,int,int8,int16,%
%      int32,int64,bool,float32,float64,complex64,complex128,byte,rune,uintptr,%
%      error,interface},%
%    morekeywords=[3]{map,slice,make,new,nil,len,cap,copy,close,true,false,%
%      delete,append,real,imag,complex,chan,},%
%    morekeywords=[4]{for,break,continue,range,go,goto,switch,case,fallthrough,if,%
%      else,default,},%
%    morekeywords=[5]{Println,Printf,Error,Print,},%
%    sensitive=true,%
%    morecomment=[l]{//},%
%    morecomment=[s]{/*}{*/},%
%    morestring=[b]',%
%    morestring=[b]",%
%    morestring=[s]{`}{`},%
% }

% \lstset{
% 	literate=
% 	{а}{{\selectfont\char224}}1
% 	{б}{{\selectfont\char225}}1
% 	{в}{{\selectfont\char226}}1
% 	{г}{{\selectfont\char227}}1
% 	{д}{{\selectfont\char228}}1
% 	{е}{{\selectfont\char229}}1
% 	{ё}{{\"e}}1
% 	{ж}{{\selectfont\char230}}1
% 	{з}{{\selectfont\char231}}1
% 	{и}{{\selectfont\char232}}1
% 	{й}{{\selectfont\char233}}1
% 	{к}{{\selectfont\char234}}1
% 	{л}{{\selectfont\char235}}1
% 	{м}{{\selectfont\char236}}1
% 	{н}{{\selectfont\char237}}1
% 	{о}{{\selectfont\char238}}1
% 	{п}{{\selectfont\char239}}1
% 	{р}{{\selectfont\char240}}1
% 	{с}{{\selectfont\char241}}1
% 	{т}{{\selectfont\char242}}1
% 	{у}{{\selectfont\char243}}1
% 	{ф}{{\selectfont\char244}}1
% 	{х}{{\selectfont\char245}}1
% 	{ц}{{\selectfont\char246}}1
% 	{ч}{{\selectfont\char247}}1
% 	{ш}{{\selectfont\char248}}1
% 	{щ}{{\selectfont\char249}}1
% 	{ъ}{{\selectfont\char250}}1
% 	{ы}{{\selectfont\char251}}1
% 	{ь}{{\selectfont\char252}}1
% 	{э}{{\selectfont\char253}}1
% 	{ю}{{\selectfont\char254}}1
% 	{я}{{\selectfont\char255}}1
% 	{А}{{\selectfont\char192}}1
% 	{Б}{{\selectfont\char193}}1
% 	{В}{{\selectfont\char194}}1
% 	{Г}{{\selectfont\char195}}1
% 	{Д}{{\selectfont\char196}}1
% 	{Е}{{\selectfont\char197}}1
% 	{Ё}{{\"E}}1
% 	{Ж}{{\selectfont\char198}}1
% 	{З}{{\selectfont\char199}}1
% 	{И}{{\selectfont\char200}}1
% 	{Й}{{\selectfont\char201}}1
% 	{К}{{\selectfont\char202}}1
% 	{Л}{{\selectfont\char203}}1
% 	{М}{{\selectfont\char204}}1
% 	{Н}{{\selectfont\char205}}1
% 	{О}{{\selectfont\char206}}1
% 	{П}{{\selectfont\char207}}1
% 	{Р}{{\selectfont\char208}}1
% 	{С}{{\selectfont\char209}}1
% 	{Т}{{\selectfont\char210}}1
% 	{У}{{\selectfont\char211}}1
% 	{Ф}{{\selectfont\char212}}1
% 	{Х}{{\selectfont\char213}}1
% 	{Ц}{{\selectfont\char214}}1
% 	{Ч}{{\selectfont\char215}}1
% 	{Ш}{{\selectfont\char216}}1
% 	{Щ}{{\selectfont\char217}}1
% 	{Ъ}{{\selectfont\char218}}1
% 	{Ы}{{\selectfont\char219}}1
% 	{Ь}{{\selectfont\char220}}1
% 	{Э}{{\selectfont\char221}}1
% 	{Ю}{{\selectfont\char222}}1
% 	{Я}{{\selectfont\char223}}1
% }

% % Работа с изображениями и таблицами; переопределение названий по ГОСТу
% \usepackage{caption}
% \captionsetup[figure]{name={Рисунок},labelsep=endash}
% % \captionsetup[table]{singlelinecheck=false, labelsep=endash}
% \captionsetup[table]{justification=centering, singlelinecheck=false, labelsep=endash}
% % \captionsetup[lstlisting]{position=top, justification=centering, singlelinecheck=false}

% \usepackage[justification=centering]{caption} % Настройка подписей float объектов	

% \usepackage{csvsimple}

% \usepackage{ulem} % Нормальное нижнее подчеркивание
% \usepackage{hhline} % Двойная горизонтальная линия в таблицах
% \usepackage[figure,table]{totalcount} % Подсчет изображений, таблиц
% \usepackage{rotating} % Поворот изображения вместе с названием
% \usepackage{lastpage} % Для подсчета числа страниц

% \makeatletter
% \renewcommand\@biblabel[1]{#1.} % [1] -> 1. in bibliography
% \makeatother

% \usepackage{float} % Place figures anywhere you want; ignore floating

% \usepackage{ragged2e} % Перенос слов на следующую строку
% \usepackage{pdfpages}

% \usepackage{threeparttable}

% \usepackage{blindtext}

% % CW-specific

% \usepackage{xparse} % \NewDocumentCommand for creating custom commands

% % \usepackage{tabularx}
% \usepackage{adjustbox} % Hyperlinks
% % \usepackage[table]{xcolor}

% \renewcommand{\underset}[2]{\ensuremath{\mathop{\mbox{#2}}\limits_{\mbox{\scriptsize #1}}}} % Use \underset in normal environment, not only in math

% % Write stuff under underlined text
% \NewDocumentCommand{\ulinetext}{O{3cm} O{c} m m} % O - optional; m - mandatory
% {\underset{#3}{\uline{\makebox[#1][#2]{#4}}}}

% % Font configuration
% % \usepackage{fontspec}
% % \setmainfont{Liberation Serif} % Liberation Serif} % ~ Times New Roman
% % \setsansfont{Roboto} % Liberation Sans} % ~ Arial
% % \setsansfont{Liberation Sans} % ~ Arial
% % \setmonofont{Liberation Mono} % ~ Courier New

% \usepackage{totcount}

% % Подсчет количества рисунков и таблиц
% \regtotcounter{figure}
% \regtotcounter{table}



\usepackage{adjustbox}
\usepackage{amsfonts}
\usepackage{amsmath}
\usepackage{array}
\usepackage[english,russian]{babel}
\usepackage{blindtext}
\usepackage{bmstu-title}
\usepackage{caption}
\usepackage{chngcntr}
\usepackage{cmap}
\usepackage{csvsimple}
\usepackage{enumitem} 
\usepackage{etoolbox}
\usepackage[14pt]{extsizes}
\usepackage{float}
\usepackage[T2A]{fontenc}
\usepackage{geometry}
\usepackage{graphicx}
\usepackage{hhline}
\usepackage[pdftex]{hyperref}
\usepackage{lastpage}     % для \pageref{LastPage}
\usepackage{totcount}     % для \total{...}
\regtotcounter{figure}    % регистрирует счётчик figure для \total
\regtotcounter{table}     % (опционально) если используете \total{table}
\usepackage{indentfirst}
\usepackage[utf8]{inputenc}
\usepackage{lastpage}
\usepackage{listings}
\usepackage{ltablex}
\usepackage{microtype} 
\usepackage{multirow}
\usepackage{pdfpages}
\usepackage{ragged2e}
\usepackage{rotating}
\usepackage{setspace}
\usepackage{siunitx}
\usepackage{tabularx}
\usepackage{tcolorbox}
\usepackage{titletoc}
\usepackage{totcount}
\usepackage{titlesec}[explicit]
\usepackage{ulem}
\usepackage{url}
\usepackage[table]{xcolor}
\usepackage{xparse}

% Настройка полей
% ============================
\geometry{left=30mm}
\geometry{right=15mm}
\geometry{top=20mm}
\geometry{bottom=20mm}
% ============================

% Форматирование секций
% ============================
\titleformat{name=\section,numberless}[block]{\normalfont\large\bfseries\centering}{}{0pt}{}
\titleformat{\section}[block]{\normalfont\large\bfseries}{\thesection}{1em}{}

\titleformat{\subsection}[hang]
{\bfseries\large}{\thesubsection}{1em}{}
\titlespacing\subsection{\parindent}{*2}{*2}

\titleformat{\subsubsection}[hang]
{\bfseries\large}{\thesubsubsection}{1em}{}
\titlespacing\subsubsection{\parindent}{*2}{*2}

\titlespacing\section{\parindent}{*4}{*4}

% Отточие у секций
\titlecontents{section}[0pt]
  {\addvspace{1em}}
  {\bfseries\thecontentslabel\quad}
  {\bfseries}
  {\normalfont\hspace{0.5em}\titlerule*[1pc]{.}\contentspage}
  []
% ============================

% Общее форматирование
% ============================
% Полуторный интервал
\onehalfspacing

% Убираем увеличенные пробелы после точек
\frenchspacing

% Красная строка
\setlength\parindent{1.25cm}
% ============================

% Сокрытие цветных рамок вокруг ссылок
\hypersetup{hidelinks}

% Настройка списков
% ============================
\setlist{
    labelsep=0.3cm,
    itemsep=0pt,
    parsep=0pt,
    topsep=0pt,
    itemindent=0pt
}

\setlist[enumerate]{
    label=\arabic*), 
    leftmargin=2cm,
}
\setlist[enumerate,2]{
    label=\arabic*), 
    leftmargin=*,
}


\setlist[itemize]{
    label=---,
    leftmargin=2cm,
    labelsep=0.3cm
}
\setlist[itemize,2]{
    label=---,
    leftmargin=*,
    labelsep=0.3cm
}
% ============================

% Настройки листингов
% ============================
\lstset{
	backgroundcolor=\color{white},
	basicstyle=\footnotesize\ttfamily\linespread{0.8},
	breakatwhitespace=true,	
	breaklines=true,
	captionpos=b,
	commentstyle=\color{green!50!black},
	escapeinside={\#*}{*)},	
	frame=single,
	inputencoding=utf8,
	keywordstyle=\color{blue},
	language=C,
	numbers=left,
	numbersep=5pt,
	numbersep=5pt,
	numberstyle=\tiny,
	showstringspaces=false,
	stringstyle=\color{red!90!black}, 
	tabsize=4,
}
\lstset{
	morekeywords={size_t, likely}
}
\lstset{
	literate=
	{а}{{\selectfont\char224}}1
	{б}{{\selectfont\char225}}1
	{в}{{\selectfont\char226}}1
	{г}{{\selectfont\char227}}1
	{д}{{\selectfont\char228}}1
	{е}{{\selectfont\char229}}1
	{ё}{{\"e}}1
	{ж}{{\selectfont\char230}}1
	{з}{{\selectfont\char231}}1
	{и}{{\selectfont\char232}}1
	{й}{{\selectfont\char233}}1
	{к}{{\selectfont\char234}}1
	{л}{{\selectfont\char235}}1
	{м}{{\selectfont\char236}}1
	{н}{{\selectfont\char237}}1
	{о}{{\selectfont\char238}}1
	{п}{{\selectfont\char239}}1
	{р}{{\selectfont\char240}}1
	{с}{{\selectfont\char241}}1
	{т}{{\selectfont\char242}}1
	{у}{{\selectfont\char243}}1
	{ф}{{\selectfont\char244}}1
	{х}{{\selectfont\char245}}1
	{ц}{{\selectfont\char246}}1
	{ч}{{\selectfont\char247}}1
	{ш}{{\selectfont\char248}}1
	{щ}{{\selectfont\char249}}1
	{ъ}{{\selectfont\char250}}1
	{ы}{{\selectfont\char251}}1
	{ь}{{\selectfont\char252}}1
	{э}{{\selectfont\char253}}1
	{ю}{{\selectfont\char254}}1
	{я}{{\selectfont\char255}}1
	{А}{{\selectfont\char192}}1
	{Б}{{\selectfont\char193}}1
	{В}{{\selectfont\char194}}1
	{Г}{{\selectfont\char195}}1
	{Д}{{\selectfont\char196}}1
	{Е}{{\selectfont\char197}}1
	{Ё}{{\"E}}1
	{Ж}{{\selectfont\char198}}1
	{З}{{\selectfont\char199}}1
	{И}{{\selectfont\char200}}1
	{Й}{{\selectfont\char201}}1
	{К}{{\selectfont\char202}}1
	{Л}{{\selectfont\char203}}1
	{М}{{\selectfont\char204}}1
	{Н}{{\selectfont\char205}}1
	{О}{{\selectfont\char206}}1
	{П}{{\selectfont\char207}}1
	{Р}{{\selectfont\char208}}1
	{С}{{\selectfont\char209}}1
	{Т}{{\selectfont\char210}}1
	{У}{{\selectfont\char211}}1
	{Ф}{{\selectfont\char212}}1
	{Х}{{\selectfont\char213}}1
	{Ц}{{\selectfont\char214}}1
	{Ч}{{\selectfont\char215}}1
	{Ш}{{\selectfont\char216}}1
	{Щ}{{\selectfont\char217}}1
	{Ъ}{{\selectfont\char218}}1
	{Ы}{{\selectfont\char219}}1
	{Ь}{{\selectfont\char220}}1
	{Э}{{\selectfont\char221}}1
	{Ю}{{\selectfont\char222}}1
	{Я}{{\selectfont\char223}}1
}
% ============================

% Настройка подписей 
% ============================
\captionsetup[figure]{
	name={Рисунок},
	justification=centering,
	labelsep=endash,
	singlelinecheck=on
}

\captionsetup[table]{
	name={Таблица},
	justification=raggedright,
	labelsep=endash,
	singlelinecheck=off
}

% \captionsetup[lstlisting]{
% 	name={Листинг},
% 	justification=centering,
% 	labelsep=endash,
% 	singlelinecheck=on
% }
% ============================

% Метки в библиографии
% ============================
\makeatletter
\renewcommand\@biblabel[1]{#1.}
\makeatother
% ============================

% Переопределение \underset для лучшего отображения нижних индексов в тексте
\renewcommand{\underset}[2]{\ensuremath{\mathop{\mbox{#2}}\limits_{\mbox{\scriptsize #1}}}}

% Подчёркнутый текст с пояснением снизу
\NewDocumentCommand{\ulinetext}{O{3cm} O{c} m m}
{\underset{#3}{\uline{\makebox[#1][#2]{#4}}}}

% Секция без номера, но с записью в оглавлении
\newcommand{\numberlesssection}[1]{\phantomsection\section*{#1}\addcontentsline{toc}{section}{#1}}

% Подсекция без номера, но с записью в оглавлении
\newcommand{\numberlesssubsection}[1]{\phantomsection\subsection*{#1}\addcontentsline{toc}{subsection}{#1}}

% Вставка картинок
\newcommand{\img}[3] {
	\begin{figure}[h!]
		\center{\includegraphics[height=#1]{img/#2}}
		\caption{#3}
		\label{img:#2}
	\end{figure}
}

% Картинки, растянутые до полей
\newcommand{\jimg}[2] {
	\begin{figure}[h!]
		\center{\includegraphics[width=\textwidth]{img/#1}}
		\caption{#2}
		\label{img:#1}
	\end{figure}
}

% Inline код
\newcommand{\code}[1]{{\small \ttfamily #1}}

% Центрирование в таблице, растянутой по ширине страницы
\newcolumntype{Y}{>{\centering\arraybackslash}X}

% Приложения
\newcounter{appsection}
\renewcommand{\theappsection}{\Asbuk{appsection}}

\newcommand{\appsection}{
	\newpage
	\stepcounter{appsection}
	\def\appname{ПРИЛОЖЕНИЕ~\theappsection}
	\phantomsection
	\section*{\appname}
	\addcontentsline{toc}{section}{\appname}

	\renewcommand{\thefigure}{\theappsection.\arabic{figure}}
	\renewcommand{\thetable}{\theappsection.\arabic{table}}
	\renewcommand{\thelstlisting}{\theappsection.\arabic{lstlisting}}
	\renewcommand{\theequation}{\theappsection.\arabic{equation}}

	\setcounter{figure}{0}
	\setcounter{table}{0}
	\setcounter{lstlisting}{0}
	\setcounter{equation}{0}
	\setcounter{section}{0}
	\setcounter{subsection}{0}
	\setcounter{subsubsection}{0}
}

% Сохранение ширины таблицы после переноса, для ltablex
\keepXColumns

% Листинги с автопереносом
\tcbuselibrary{listings, breakable, skins}
\newtcblisting[use counter=lstlisting]{listing}[3][]{
	boxrule=0.5pt,
	breakable,
	colback=white,
	colbacktitle=white,
	colframe=black,
	fonttitle=\color{black},
	listing only,
	listing options={
		language=#3,
		frame=none,
		xleftmargin=0pt,
		title={#1}
	},
	sharp corners,
	title={Листинг~\thetcbcounter~--~#1},
	title after break={Листинг~\thetcbcounter\ (продолжение)},
	label={#2},
}

% Листинг из файла
\newtcbinputlisting[use counter=lstlisting]{\flisting}[4][]{
	boxrule=0.5pt,
	breakable,
	colback=white,
	colbacktitle=white,
	colframe=black,
	fonttitle=\color{black},
	listing only,
	listing file={#3},
	listing options={
		language=#4,
		frame=none,
		xleftmargin=10pt
	},
	sharp corners,
	title={Листинг~\thetcbcounter~--~#1},
	title after break={Листинг~\thetcbcounter\ (продолжение)},
	label={#2}
}

% Вставка страницы A3
\newcommand{\AThreeStart}{
	\clearpage
	\pdfpagewidth=420mm
	\pdfpageheight=297mm
	\paperwidth=420mm
	\paperheight=297mm

	\textwidth=375mm
	\textheight=257mm
	\columnwidth=\textwidth
	\hsize=\textwidth
	\vsize=\textheight
}

\newcommand{\AThreeEnd}{
	\clearpage
	\pdfpagewidth=210mm
	\pdfpageheight=297mm
	\paperwidth=210mm
	\paperheight=297mm
}

\newenvironment{a3page}
	{\AThreeStart}
	{\AThreeEnd}

% Подсчет общего количества рисунков, таблиц и источников
\newcounter{figures}
\regtotcounter{figures}
\AtBeginEnvironment{figure}{\stepcounter{figures}}

\newcounter{tables}
\regtotcounter{tables}
\AtBeginEnvironment{table}{\stepcounter{tables}}

\newcounter{references}
\regtotcounter{references}
\makeatletter
\pretocmd{\bibitem}{\stepcounter{references}}{}{}
\makeatother